% ---------------------------------------------------------------------------
% Einleitung (Merkblatt 3.2: Problem-/Fragestellung, Zielsetzung, Aufbau;
% Umfang max. ~15% der Seiten)
% ---------------------------------------------------------------------------
\chapter{Einleitung}
\label{cha:einleitung}

Deutschland erlebte in den Jahren 2015 und 2016 die größte Flüchtlingswelle der jüngeren Vergangenheit \parencite{brucker2016vorlaufige}. Insbesondere Kriege, Bedrohungen der Meinungsfreiheit waren einige der Gründe, die Flüchtlinge dazu veranlassten, nach Deutschland zu fliehen \parencite[4]{brucker202510}. Der 2015 ausgebrochene Bürgerkrieg in Syrien sowie die politischen und wirtschaftlichen Probleme in Ländern wie dem Irak, Afghanistan und Eritrea waren die Hauptursachen für die Flüchtlingswelle der Jahre 2015–2016 \parencites[4]{unknown}[23--24]{brucker2016vorlaufige}. Insbesondere in den Jahren 2015 und 2016 erreichte die Zahl der nach Deutschland geflüchteten Menschen ihren Höchststand. Im Jahr 2016 stellten 722.370 Menschen einen Asylantrag in Deutschland, während diese Zahl im Jahr 2017 auf 198.317 zurückging und seitdem stetig abgenommen hat \parencite[103]{fur2019migrationsbericht}. Obwohl diese Anzahl in den letzten Jahren zurückgegangen ist, hat sie seit dieser Zeit die Politik Deutschlands in den Bereichen Asyl, Migration und Integration maßgeblich beeinflusst.

Der Erwerb der deutschen Sprache ist einer der ersten Schritte im Rahmen der Integration von Geflüchteten in den Arbeitsmarkt \parencite{Bahr2019Arbeitsmarktintegration}. Angesichts der Anforderungen vieler Berufsgruppen in Deutschland, in denen für mindestens die Hälfte der Stellen Deutschkenntnisse auf dem Niveau B2 oder höher vorausgesetzt werden, sind die für Geflüchtete angebotenen Deutschkurse von entscheidender Bedeutung für eine erfolgreiche Integration in den Arbeitsmarkt \parencite{repec:iab:iabkbe:202607}. Darüber hinaus wird die Teilnahme von Migrant:innen, deren Deutschkenntnisse unter dem Niveau B2 liegen, an einem Deutschkurs als ein zu berücksichtigender Faktor angesehen. Nicht nur Unternehmen, sondern auch 52 Prozent der Arbeitssuchenden sind der Einschätzung, dass mangelnde Deutschkenntnisse bei der Arbeitssuche ein mögliches Hindernis darstellen könnten \parencite{repec:iab:iabkbe:202607}.

Die Geflüchteten verfügten anfangs nur über geringe Deutschkenntnisse. Laut Daten aus dem Jahr 2016 gaben lediglich 16 Prozent der Geflüchteten an, gut Deutsch sprechen zu können \parencite[7]{brucker2018iab}. Bei etwa der Hälfte der Geflüchteten wurde berichtet, dass sie „nicht oder eher schlecht Deutsch sprechen“ \parencite[7]{brucker2018iab}.

Bemerkenswert ist auch die anfänglich niedrige Beschäftigungsquote der Geflüchteten in Deutschland. Während im Jahr 2015 9,9 Prozent der nach Deutschland gekommenen Geflüchteten erwerbstätig waren, lag diese Quote im Jahr 2016 bei 6,2 Prozent \parencite{brucker2017arbeitsmarktintegration}. Im Laufe der Zeit war jedoch ein Anstieg der Beschäftigungsquote der Geflüchteten zu beobachten \parencite{brucker2024institutionelle}.

Nach Angaben aus dem Jahr 2016 waren lediglich 3 Prozent der Geflüchteten ohne jegliche Deutschkenntnisse erwerbstätig, während von den Geflüchteten mit geringen Deutschkenntnissen lediglich 7 Prozent erwerbstätig waren (blau) \parencite[9]{sirries2017arbeitsmarktintegration}. Im Vergleich zu anderen Möglichkeiten des Sprachenlernens profitieren Geflüchtete besonders von strukturierten Lernangeboten wie Sprachkursen \parencite{kristen2022deutschkenntnisse}. Aus diesem Grund gewinnt der Einfluss des lokalen Sprachkursangebots auf die Beschäftigung gesellschaftliche Relevanz.

Der Aufenthaltsstatus gilt als wichtiger Faktor für den Zugang von Geflüchteten zum Arbeitsmarkt. Während rechtliche Hindernisse für den Zugang von Geflüchteten mit anerkanntem Aufenthaltsstatus zum Arbeitsmarkt aufgehoben werden, bestehen für Geflüchtete, deren Aufenthaltsstatus noch nicht geklärt ist oder deren Antrag abgelehnt wurde, weiterhin bürokratische Hindernisse \parencite[128--129]{kosyakova2020role}. Darüber hinaus kann die ungewisse Aufenthaltssituation von Geflüchteten, deren Aufenthaltsstatus noch nicht feststeht oder deren Antrag abgelehnt wurde, die Entscheidungen von Arbeitgebenden beeinflussen \parencite{kosyakova2020role}. \textcite{brucker2024institutionelle} stellten fest, dass die Wahrscheinlichkeit einer Beschäftigung bei Personen mit anerkanntem Aufenthaltsstatus im Vergleich zu Personen mit abgelehntem oder ungewissem Status höher ist.

Darüber hinaus spielt der Aufenthaltsstatus eine zentrale Rolle beim Zugang zu Deutschkursen. Personen mit anerkanntem Status erhalten uneingeschränkten Zugang zu öffentlich geförderten Deutschkursen, während Personen, deren Asylantrag abgelehnt wurde, unter bestimmten Umständen im Rahmen einer Duldung das Recht auf Teilnahme an Sprachkursen erhalten \parencite[129]{kosyakova2020role}. Personen, deren Aufenthaltsstatus noch nicht geklärt ist, haben hingegen keinen uneingeschränkten Zugang zu öffentlich geförderten Sprachkursen \parencite[12]{kosyakova2020role}.

Der Aufenthaltsstatus wirkt sich sowohl auf die Beschäftigungschancen von Geflüchteten als auch auf deren Zugang zu staatlich geförderten Deutschkursen aus. Aus diesem Grund ist es aus soziologischer Sicht von Interesse, welche Gruppe von Geflüchteten im Hinblick auf die Integration durch Sprachkurse in Bezug auf die Beschäftigung am stärksten profitiert.


\textcite{Kanas:2023} untersuchten den kausalen Einfluss des Angebots an lokalen Sprachkursen in Deutschland auf die Integration von Flüchtlingen in den Arbeitsmarkt. Dabei wurde der Einfluss des Sprachkursangebots in den Landkreisen auf die Beschäftigungschancen von Geflüchteten untersucht. \textcite{Kanas:2023} stellten fest, dass die Beschäftigungschancen von Geflüchteten in den Landkreisen, in denen mehr lokale Sprachkurse angeboten werden, höher sind. Diese Arbeit zielt darauf ab, die empirischen Analysen aus dem Artikel von \textcite{Kanas:2023} zu reproduzieren, um die Ergebnisse dieser Arbeit erneut zu analysieren und um eine Erweiterung vorzunehmen. Der Beitrag dieser Arbeit besteht darin, zu untersuchen, ob Flüchtlinge, deren Aufenthaltsstatus anerkannt ist, im Hinblick auf ihren Erfolg auf dem Arbeitsmarkt stärker von dem Angebot an lokalen Sprachkursen in Deutschland profitieren. Anhand des Spracherwerbsmodells von \textcite{Esser:2006} wird in dieser Arbeit die Investition von Geflüchteten in Sprachkurse in Abhängigkeit von ihrem Aufenthaltsstatus erklärt und anschließend der Aufenthaltsstatus als Interaktionseffekt getestet. Demzufolge beantwortet diese Arbeit die folgende Forschungsfrage: „Profitieren Geflüchtete mit anerkanntem Asylantrag stärker von lokalen Sprachkursen im Hinblick auf ihren Arbeitsmarkterfolg?“


Diese Arbeit ist wie folgt aufgebaut: Der Forschungsstand gliedert sich in drei verschiedene Abschnitte. Zunächst wird der Zusammenhang zwischen dem Spracherwerb der Zielsprache und dem Erfolg auf dem Arbeitsmarkt erläutert. Anschließend werden Beiträge vorgestellt, die den Einfluss von Sprachkursen auf die Integration in den Arbeitsmarkt beleuchten. Im Hinblick auf die Erweiterung dieser Arbeit werden im letzten Abschnitt Studien vorgestellt, die sich mit dem Zusammenhang zwischen dem Aufenthaltsstatus und dem Spracherwerb im Hinblick auf die Integration in den Arbeitsmarkt befassen. Im letzten Abschnitt des Forschungsstands wird schließlich die Entwicklung von Integrationskursen und Teilnahmebescheinigungen auf Kreisebene dargestellt. „Theoretischer Rahmen“ wird in zwei Unterkapiteln behandelt. Zunächst werden die theoretischen Ansätze erläutert, mit denen Kanas und Kosyakova die Auswirkungen des lokalen Kursangebots auf die Beschäftigung erklären. Der zweite Teil des „Theoretischen Teils“ basiert auf einer theoretischen Erweiterung. Anhand des Sprachlernmodells von Esser (2006) wird theoretisch erläutert, ob sich die Auswirkungen des lokalen Kursangebots auf die Beschäftigung je nach Asylstatus unterscheiden. Im methodischen Teil werden die sowohl in „Reproduction“ als auch in „Erweiterung“ verwendeten Datensätze, die Replikationsstrategie, die Operationalisierung der Variablen sowie die Bildung der Stichprobe erläutert. Im Ergebnisteil wird die Forschungsfrage anhand der Ergebnisse von „Reproduction“ und „Erweiterung“ beantwortet.
